% GENERATED FILE. Edit canonical lessons, then run npm run generate:books.
% canonical-source-hash: fnv1a64:9e4807dee61a741a
% canonical-lessons: HI-C296-raftar, HI-C296-takrana, HI-C296-savari, HI-C296-pahiya, HI-C296-palatna

\chapter{A Crash on the Road}
\label{ch:hi-a-crash-on-the-road}
\hlchaptermodality{296}

\begin{chapteropening}
\textbf{By the end of this chapter:} \emph{I can say how an accident on the road happened.}

Complete the last lesson of chapter 296: I can say how an accident on the road happened.
\end{chapteropening}

\section[\texorpdfstring{\hi{रफ़्तार} (raftār)}{रफ़्तार (raftār)}]{\hi{रफ़्तार} (raftār) — speed}

\label{lesson:HI-C296-raftar}



\begin{quote}
Before the new one: say the Hindi for a fine, then the Hindi for handcuffs.
\end{quote}

\subsection*{You'll want to know: \hi{रफ़्तार}}
\textbf{\hi{रफ़्तार}} — \emph{raftār} — \textquotedblleft{}speed\textquotedblright{}.

\begin{quote}
\hi{ट्रक} \hi{की} \hi{रफ़्तार} \hi{बस} \hi{से} \hi{ज़्यादा} \hi{थी।}
\end{quote}

\begin{quote}
\emph{ṭrak kī raftār bas se zyādā thī.}
\end{quote}

\begin{quote}
The lorry was going faster than the bus.
\end{quote}

Feminine.

\begin{cousinweb}
\textbf{\hi{रफ़्तार}} (\emph{raftār}) is usually traced to Persian \emph{raftār}, \textquotedblleft{}a going, a gait\textquotedblright{}, from the Persian verb to go. It is feminine: \textbf{\hi{तेज़} \hi{रफ़्तार}}, high speed.
\end{cousinweb}

\begin{grammarlens}[title={What you've built}]
One more word for telling what happened.
\end{grammarlens}

\subsection*{Your turn}
\begin{itemize}
\raggedright
  \item \emph{Say it:} \emph{raftār}
  \item \emph{Say it:} \emph{raftār}, once more
  \item \emph{Say it:} \emph{ṭrak kī raftār bas se zyādā thī.} — The lorry was going faster than the bus.
  \item \emph{Recall:} say the Hindi for a fine, then the Hindi for handcuffs, then say \emph{raftār} again
\end{itemize}

\begin{tcolorbox}[breakable,colback=teal!4,colframe=teal!35!black,title={Before you move on}]
What does \hi{रफ़्तार} mean? (\textquotedblleft{}Speed\textquotedblright{}.) Say it once more.
\end{tcolorbox}
\section[\texorpdfstring{\hi{टकराना} (ṭakrānā)}{टकराना (ṭakrānā)}]{\hi{टकराना} (ṭakrānā) — to collide; to crash into}

\label{lesson:HI-C296-takrana}



\begin{quote}
Before the new one: say the Hindi for handcuffs, then the Hindi for speed.
\end{quote}

\subsection*{You'll want to know: \hi{टकराना}}
\textbf{\hi{टकराना}} — \emph{ṭakrānā} — \textquotedblleft{}to collide; to crash into\textquotedblright{}.

\begin{quote}
\hi{ट्रक} \hi{एक} \hi{बस} \hi{से} \hi{टकराया।}
\end{quote}

\begin{quote}
\emph{ṭrak ek bas se ṭakrāyā.}
\end{quote}

\begin{quote}
The lorry crashed into a bus.
\end{quote}

Past (it …): \hi{वह} \hi{ट्रक} \hi{पेड़} \hi{से} \hi{टकराया}, \emph{vah ṭrak peṛ se ṭakrāyā}.

\begin{cousinweb}
\textbf{\hi{टकराना}} (\emph{ṭakrānā}) is the verb made on \emph{ṭakkar}, a blow or a crash, which CDIAL 5424 lists for Hindi under \emph{ṭakkarā}, \textquotedblleft{}a blow on the head\textquotedblright{}. What you hit takes \textbf{\hi{से}}: \textbf{\hi{बस} \hi{से} \hi{टकराया}}.
\end{cousinweb}

\begin{grammarlens}[title={What you've built}]
One more word for telling what happened.
\end{grammarlens}

\subsection*{Your turn}
\begin{itemize}
\raggedright
  \item \emph{Say it:} \emph{ṭakrānā}
  \item \emph{Say it:} \emph{ṭakrānā}, once more
  \item \emph{Say it:} \emph{ṭrak ek bas se ṭakrāyā.} — The lorry crashed into a bus.
  \item \emph{Recall:} say the Hindi for handcuffs, then the Hindi for speed, then say \emph{ṭakrānā} again
\end{itemize}

\begin{tcolorbox}[breakable,colback=teal!4,colframe=teal!35!black,title={Before you move on}]
What does \hi{टकराना} mean? (\textquotedblleft{}To collide; to crash into\textquotedblright{}.) Say it once more.
\end{tcolorbox}
\section[\texorpdfstring{\hi{सवारी} (savārī)}{सवारी (savārī)}]{\hi{सवारी} (savārī) — a passenger; a ride}

\label{lesson:HI-C296-savari}



\begin{quote}
Before the new one: say the Hindi for speed, then the Hindi for to collide.
\end{quote}

\subsection*{You'll want to know: \hi{सवारी}}
\textbf{\hi{सवारी}} — \emph{savārī} — \textquotedblleft{}a passenger; a ride\textquotedblright{}.

\begin{quote}
\hi{बस} \hi{में} \hi{चालीस} \hi{सवारियाँ} \hi{थीं।}
\end{quote}

\begin{quote}
\emph{bas meṁ cālīs savāriyā̃ thīṁ.}
\end{quote}

\begin{quote}
There were forty passengers on the bus.
\end{quote}

Feminine.

\begin{cousinweb}
\textbf{\hi{सवारी}} (\emph{savārī}) is usually traced to Persian \emph{savārī}, \textquotedblleft{}riding\textquotedblright{}, from \emph{savār}, a rider. It is feminine whatever the passenger: \textbf{\hi{चालीस} \hi{सवारियाँ}}.
\end{cousinweb}

\begin{grammarlens}[title={What you've built}]
One more word for telling what happened.
\end{grammarlens}

\subsection*{Your turn}
\begin{itemize}
\raggedright
  \item \emph{Say it:} \emph{savārī}
  \item \emph{Say it:} \emph{savārī}, once more
  \item \emph{Say it:} \emph{bas meṁ cālīs savāriyā̃ thīṁ.} — There were forty passengers on the bus.
  \item \emph{Recall:} say the Hindi for speed, then the Hindi for to collide, then say \emph{savārī} again
\end{itemize}

\begin{tcolorbox}[breakable,colback=teal!4,colframe=teal!35!black,title={Before you move on}]
What does \hi{सवारी} mean? (\textquotedblleft{}A passenger; a ride\textquotedblright{}.) Say it once more.
\end{tcolorbox}
\section[\texorpdfstring{\hi{पहिया} (pahiyā)}{पहिया (pahiyā)}]{\hi{पहिया} (pahiyā) — a wheel}

\label{lesson:HI-C296-pahiya}



\begin{quote}
Before the new one: say the Hindi for to collide, then the Hindi for a passenger.
\end{quote}

\subsection*{You'll want to know: \hi{पहिया}}
\textbf{\hi{पहिया}} — \emph{pahiyā} — \textquotedblleft{}a wheel\textquotedblright{}.

\begin{quote}
\hi{स्कूटर} \hi{का} \hi{पिछला} \hi{पहिया} \hi{टूट} \hi{गया।}
\end{quote}

\begin{quote}
\emph{skūṭar kā pichlā pahiyā ṭūṭ gayā.}
\end{quote}

\begin{quote}
The scooter's back wheel broke.
\end{quote}

Masculine.

\begin{cousinweb}
\textbf{\hi{पहिया}} (\emph{pahiyā}) goes back to Sanskrit \emph{pradhi}, \textquotedblleft{}the rim of a wheel\textquotedblright{} (CDIAL 8672), where CDIAL records Hindi \emph{pahiyā}. Its plural is \textbf{\hi{पहिए}}.
\end{cousinweb}

\begin{grammarlens}[title={What you've built}]
One more word for telling what happened.
\end{grammarlens}

\subsection*{Your turn}
\begin{itemize}
\raggedright
  \item \emph{Say it:} \emph{pahiyā}
  \item \emph{Say it:} \emph{pahiyā}, once more
  \item \emph{Say it:} \emph{skūṭar kā pichlā pahiyā ṭūṭ gayā.} — The scooter's back wheel broke.
  \item \emph{Recall:} say the Hindi for to collide, then the Hindi for a passenger, then say \emph{pahiyā} again
\end{itemize}

\begin{tcolorbox}[breakable,colback=teal!4,colframe=teal!35!black,title={Before you move on}]
What does \hi{पहिया} mean? (\textquotedblleft{}A wheel\textquotedblright{}.) Say it once more.
\end{tcolorbox}
\section[\texorpdfstring{\hi{पलटना} (palaṭnā)}{पलटना (palaṭnā)}]{\hi{पलटना} (palaṭnā) — to overturn}

\label{lesson:HI-C296-palatna}



\begin{quote}
Before the new one: say the Hindi for a passenger, then the Hindi for a wheel.
\end{quote}

\subsection*{You'll want to know: \hi{पलटना}}
\textbf{\hi{पलटना}} — \emph{palaṭnā} — \textquotedblleft{}to overturn\textquotedblright{}.

\begin{quote}
\hi{मोड़} \hi{पर} \hi{बस} \hi{पलट} \hi{गई।}
\end{quote}

\begin{quote}
\emph{moṛ par bas palaṭ gaī.}
\end{quote}

\begin{quote}
The bus overturned on the bend.
\end{quote}

Past (it …): \hi{वह} \hi{गाड़ी} \hi{पलट} \hi{गई}, \emph{vah gāṛī palaṭ gaī}.

\begin{cousinweb}
is inherited from \emph{*pallaṭṭ}, a form CDIAL reconstructs (CDIAL 7968). It is a whole turning over: a bus \textbf{\hi{पलट} \hi{गई}}, and a page or a decision can \textbf{\hi{पलटना}} too.
\end{cousinweb}

\begin{grammarlens}[title={What you've built}]
One more word for telling what happened.
\end{grammarlens}

\subsection*{Your turn}
\begin{itemize}
\raggedright
  \item \emph{Say it:} \emph{palaṭnā}
  \item \emph{Say it:} \emph{palaṭnā}, once more
  \item \emph{Say it:} \emph{moṛ par bas palaṭ gaī.} — The bus overturned on the bend.
  \item \emph{Recall:} say the Hindi for a passenger, then the Hindi for a wheel, then say \emph{palaṭnā} again
\end{itemize}

\begin{tcolorbox}[breakable,colback=teal!4,colframe=teal!35!black,title={Before you move on}]
What does \hi{पलटना} mean? (\textquotedblleft{}To overturn\textquotedblright{}.) Say it once more.
\end{tcolorbox}
